% ════════════════════════════════════════════════════════════════════════════════════════
\section{Part 2: Manipulating and normalizing the data}

\begin{frame}{Part 2 roadmap: 40 min, in the notebook}
  \small
  \begin{tabular}{@{}p{0.07\linewidth}p{0.6\linewidth}p{0.2\linewidth}@{}}
    \toprule
    \textbf{Step} & \textbf{Question} & \textbf{Time}\\\midrule
    \textbf{2.1} & Prepare the object & 3 min\\[3pt]
    \textbf{2.2} & What does the data look like? From \textbf{1 number} to the \textbf{whole matrix} & 10 min\\[3pt]
    \textbf{2.3} & What is a \textbf{count}? Why do libraries have different sizes? & 7 min\\[3pt]
    \textbf{2.4} & \textbf{Normalization}: what, why, how & 10 min\\[3pt]
    \textbf{2.5} & Why do we \textbf{filter} lowly expressed genes? & 5 min\\[3pt]
    \textbf{2.6} & \textbf{Heatmaps} before and after normalizing & 5 min\\
    \bottomrule
  \end{tabular}
  \vspace{0.3cm}

  \begin{block}{Goal}
    Understand what is inside the data and turn \alert{raw counts} into numbers we can \alert{compare} between samples and between genes.
  \end{block}
\end{frame}

\begin{frame}{2.2 From one number to the whole matrix}
  \centering
  \begin{tikzpicture}[font=\scriptsize]
    % the big matrix
    \draw[thick,fill=softgrey] (0,0) rectangle (6.2,-3.2);
    \node[align=center] at (3.4,-1.9) {\textbf{the whole matrix}\\ \numnsamples\ samples $\times$ \numngenesraw\ genes\\ \footnotesize 76 million numbers};
    \node[above,font=\tiny,text=greyink] at (3.1,0) {genes $\rightarrow$};
    \node[rotate=90,above,font=\tiny,text=greyink] at (0,-1.6) {samples $\rightarrow$};
    % the three zoom levels, at the corner
    \fill[accentred] (0,0) rectangle ++(0.2,-0.2);
    \fill[bluemid]   (0,-0.3) rectangle ++(1.0,-0.2);
    \fill[orangeslot] (0,-0.65) rectangle ++(1.0,-1.0);
    \foreach \i in {1,...,4} { \draw[white] (\i*0.2,-0.65) -- (\i*0.2,-1.65); \draw[white] (0,-0.65-\i*0.2) -- (1.0,-0.65-\i*0.2); }
    % explanations
    \node[anchor=west,draw=accentred,very thick,rounded corners,fill=accentred!6,minimum width=6.0cm,align=left,inner sep=3pt] (a) at (7.3,0.1)
      {\textbf{$1\times1$} one gene in one sample\\ \ttfamily\scriptsize adata[[sample], ["STAT1"]]};
    \node[anchor=west,draw=bluemid,very thick,rounded corners,fill=bluemid!6,minimum width=6.0cm,align=left,inner sep=3pt] (b) at (7.3,-1.15)
      {\textbf{$1\times5$} one sample, five genes\\ \ttfamily\scriptsize adata[[sample], five\_genes]};
    \node[anchor=west,draw=orangeslot,very thick,rounded corners,fill=orangeslot!6,minimum width=6.0cm,align=left,inner sep=3pt] (c) at (7.3,-2.4)
      {\textbf{$5\times5$} five samples, five genes\\ \ttfamily\scriptsize adata[five\_samples, five\_genes]};
  \end{tikzpicture}

  \vspace{0.2cm}
  \begin{itemize}\small\setlength\itemsep{3pt}
    \item Every piece comes with its \alert{metadata}: \code{.obs} describes the samples (rows), \code{.var} the genes (columns).
    \item The 5\,$\times$\,5 table is just the top-left corner of the whole matrix: what we do next applies to \emph{all} of it.
  \end{itemize}
\end{frame}

\begin{frame}{2.3 What is a count?}
  \centering
  \begin{tikzpicture}[font=\scriptsize, every node/.style={align=center}]
    % molecules
    \node[font=\scriptsize\bfseries] at (1.2,1.15) {3 mRNA molecules of\\ one gene (all cells summed)};
    \foreach \i/\col/\u in {0/accentred/AAGT, 1/bluemid/CCTA, 2/greenslot/GGAC} {
      \draw[thick,\col] (\i*0.95,0.1) .. controls (\i*0.95+0.2,0.4) and (\i*0.95+0.4,-0.2) .. (\i*0.95+0.7,0.1);
      \node[fill=\col,text=white,inner sep=1.5pt,font=\tiny\ttfamily] at (\i*0.95+0.35,-0.28) {\u};
    }
    \draw[-{Stealth},very thick] (3.2,0) -- node[above,font=\tiny] {amplification\\ + sequencing} (5.2,0);
    % reads
    \node[font=\scriptsize\bfseries] at (7.0,1.15) {6 reads};
    \foreach \i/\col/\u in {0/accentred/AAGT, 1/accentred/AAGT, 2/bluemid/CCTA, 3/greenslot/GGAC, 4/greenslot/GGAC, 5/greenslot/GGAC} {
      \node[fill=\col,text=white,inner sep=2pt,font=\tiny\ttfamily] at (5.6+\i*0.62,0) {\u};
    }
    \draw[-{Stealth},very thick] (9.6,0) -- node[above,font=\tiny] {collapse reads\\ with the same UMI} (11.6,0);
    \node[draw=accentred,rounded corners,fill=accentred!8,inner sep=4pt,font=\small\bfseries] at (12.9,0) {count = 3};
  \end{tikzpicture}
  \vspace{0.3cm}

  \begin{columns}[T]
    \begin{column}{0.52\linewidth}\small
      \begin{itemize}\setlength\itemsep{3pt}
        \item A \textbf{UMI} (\emph{unique molecular identifier}) is a random tag added to each \emph{original} molecule.
        \item Counting \alert{distinct UMIs} instead of reads removes the PCR duplicates: one UMI $=$ one molecule.
      \end{itemize}
    \end{column}
    \begin{column}{0.44\linewidth}\small
      \begin{itemize}\setlength\itemsep{3pt}
        \item In a pseudobulk the count is \alert{summed over all the cells} of the sample.
        \item So a count is a non-negative integer that depends on the \textbf{biology} and on \textbf{how much material} we have.
      \end{itemize}
    \end{column}
  \end{columns}
\end{frame}

\begin{frame}{2.3 Library size: cells $\times$ UMIs per cell}
  \centering
  \figfit[0.5]{p2_libsize}\\[-1pt]
  \begin{columns}[c]
    \begin{column}{0.4\linewidth}
      \centering\scriptsize
      \begin{tabular}{lrrr}\toprule
 & cells & UMIs / cell & library size\\\midrule
CD4 T & 3\,041 & 2\,989 & 8.5\,M\\
CD8 T & 1\,302 & 3\,104 & 3.8\,M\\
NK & 407 & 3\,147 & 1.2\,M\\
Monocytes & 1\,238 & 6\,747 & 7.6\,M\\
\bottomrule\end{tabular}

    \end{column}
    \begin{column}{0.58\linewidth}\footnotesize
      \begin{itemize}\setlength\itemsep{2pt}
        \item \textbf{Library size} $L_i=\sum_j\text{count}_{ij}$: total counts of sample $i$.
        \item CD4 T cells are the most abundant, NK cells the rarest ($\sim$400 cells).
        \item Monocytes have \alert{\numumiMono\ UMIs per cell} vs $\sim$\numumiNK\ in lymphocytes: bigger cells, more mRNA.
        \item PBS controls are the sum of 6 wells: \alert{\numpbsratio$\times$ deeper}.
      \end{itemize}
    \end{column}
  \end{columns}
\end{frame}

\begin{frame}{2.4 Why normalize? An example}
  \begin{columns}[c]
    \begin{column}{0.56\linewidth}
      \figfit[0.62]{p2_toy}
    \end{column}
    \begin{column}{0.42\linewidth}\small
      \textit{Is ``cytokine'' used more in a 10-page article or in a 100-page review?} Compare \alert{frequencies}, not counts.
      \vspace{0.2cm}
      \begin{itemize}\setlength\itemsep{3pt}
        \item Same donor: a PBS control (deep) and an IFN-$\gamma$ sample (5$\times$ shallower). Gene B is an interferon-response gene.
        \item \alert{Raw counts}: every gene looks \emph{lower} with IFN-$\gamma$, even the induced interferon-response gene B.
        \item \alert{CPM}: A, C, D are unchanged and B is $4\times$ up. The right answer.
      \end{itemize}
    \end{column}
  \end{columns}
\end{frame}

\begin{frame}{CPM: counts per million}
  \[
    \text{CPM}_{ij}\;=\;\frac{\text{count}_{ij}}{L_i}\times 10^{6},\qquad L_i=\sum_{j}\text{count}_{ij}
  \]
  \centering
  \figfit[0.46]{p2_actb}\\[-2pt]
  \small \gene{ACTB}, a housekeeping gene, in monocytes (PBS samples included). After CPM it no longer follows the library size.
\end{frame}

\begin{frame}{The logarithm: across genes}
  \begin{columns}[c]
    \begin{column}{0.6\linewidth}
      \figfit[0.6]{p2_logvar}
    \end{column}
    \begin{column}{0.38\linewidth}\small
      \begin{itemize}\setlength\itemsep{3pt}
        \item In CPM, \alert{variance rises steeply with the mean}.
        \item The 20 most variable genes hold \numtoptwentycpm\ of the total variance in CPM\ldots{}
        \item \ldots{}and only \numtoptwentylog\ after $\log_2(\text{CPM}+1)$.
        \item PCA, clustering and $t$-tests are driven by variance: without the log, a few abundant genes decide everything.
      \end{itemize}
    \end{column}
  \end{columns}
\end{frame}

\begin{frame}{The logarithm and the pseudocount}
  \[ \text{log-CPM}_{ij}=\log_2\bigl(\text{CPM}_{ij}+1\bigr) \]
  \begin{columns}[T]
    \begin{column}{0.45\linewidth}
      \footnotesize\textbf{Fold changes become symmetric}\\[3pt]
      \scriptsize\setlength\tabcolsep{2.5pt}
      \begin{tabular}{@{}lccccccc@{}}
        \toprule
        fold change & $\times4$ & $\times2$ & $\times1.4$ & $\times1$ & $\times0.7$ & $\times0.5$ & $\times0.25$\\
        $\lfc$ & $+2$ & $+1$ & $+0.5$ & $0$ & $-0.5$ & $-1$ & $-2$\\
        \bottomrule
      \end{tabular}\\[4pt]
      The difference of two log-values \emph{is} the log fold change.
    \end{column}
    \begin{column}{0.53\linewidth}
      \footnotesize\textbf{Why ``+1''?}\\[3pt]
      \setlength\tabcolsep{2.5pt}
      \begin{tabular}{@{}lrrrr@{}}
        \toprule
        gene & ctrl & trt & $\lfc$ & $\lfc^{+1}$\\\midrule
        almost silent & 0.1 & 0.4 & $+2.00$ & $+0.35$\\
        silent $\to$ detected & 0 & 0.5 & $+\infty$ & $+0.58$\\
        moderate & 20 & 40 & $+1.00$ & $+0.97$\\
        high & 500 & 1000 & $+1.00$ & $+1.00$\\
        \bottomrule
      \end{tabular}\\[2pt]
      \scriptsize ctrl, trt = CPM; $\lfc^{+1}$ uses CPM\,$+1$.\\[3pt]
      The pseudocount shrinks fold changes of tiny numbers: ``I don't trust ratios of one or two molecules''.
    \end{column}
  \end{columns}
\end{frame}

\begin{frame}{2.5 Why filter lowly expressed genes?}
  \centering
  \figfit[0.38]{p2_poisson}\\[-2pt]
  \small The \alert{same gene measured twice} (no true change), at four expression levels (Poisson noise only)
  \begin{itemize}\small\setlength\itemsep{3pt}
    \item Low counts are \alert{pure sampling noise}: with 2 vs 4 molecules, a ``doubling'' is bad luck.
    \item They give \alert{false fold changes} and \alert{more tests} to correct for (Part 4).
  \end{itemize}
\end{frame}

\begin{frame}{The rule: $\geq$ 10 counts in $\geq$ 12 samples}
  \begin{columns}[c]
    \begin{column}{0.56\linewidth}
      \figfit[0.6]{p2_filter}
    \end{column}
    \begin{column}{0.42\linewidth}\small
      \begin{itemize}\setlength\itemsep{4pt}
        \item 12 = the number of donors: the size of one experimental group (the idea of \texttt{filterByExpr} in edgeR).
        \item A gene that never reaches 10 counts in a whole group cannot be tested anywhere.
        \item \alert{\numngenes\ of \numngenesraw\ genes} pass. The \numngenesremoved\ removed ones carry 0.011\,\% of all counts (e.g.\ olfactory receptors).
      \end{itemize}
    \end{column}
  \end{columns}
\end{frame}

\begin{frame}{2.6 Heatmaps: cell-type markers, before and after}
  \centering
  \figfit[0.72]{p2_heat_markers}
  \vspace{-0.15cm}\\
  \scriptsize
  \textbf{z-score per gene across samples:} $z_{gi}=(x_{gi}-\bar x_g)/s_g$ (red: above the gene's average). Columns: mean of the 12 donors. \alert{Before}: the PBS columns (deepest libraries) glow. \alert{After}: clean cell-type blocks.
\end{frame}

\begin{frame}{2.6 Heatmaps: genes of the immune response}
  \begin{columns}[c]
    \begin{column}{0.71\linewidth}
      \figfit[0.8]{p2_heat_immune}
    \end{column}
    \begin{column}{0.27\linewidth}\scriptsize
      Four cytokines, four pathways:
      \begin{itemize}\setlength\itemsep{3pt}
        \item \textbf{IFN-$\gamma$ / STAT1}: \gene{GBP1}, \gene{IDO1}, \gene{CXCL9}\ldots, strongest in monocytes
        \item \textbf{IL-4 / STAT6}: \gene{ALOX15}, \gene{CCL17}, \gene{FCER2}, monocytes
        \item \textbf{IL-1$\beta$ / NF-$\kappa$B}: \gene{IL6}, \gene{CXCL8}, \gene{CCL2}, monocytes
        \item \textbf{IL-2 / STAT5}: \gene{IL2RA}, \gene{GZMB}, \gene{CISH}, NK and T cells
      \end{itemize}
      \alert{Normalized: biology. Raw: sequencing depth.}\\[2pt]
      z-scores are for visualization only.
    \end{column}
  \end{columns}
\end{frame}

\begin{frame}{Your turn in the notebook (Part 2)}
  \small
  Cells that start with \code{\# CHANGE ME} contain parameters you can modify. Change them, re-run and see what happens:
  \vspace{0.2cm}

  \begin{tabular}{@{}>{\ttfamily}p{0.27\linewidth}>{\raggedright\arraybackslash}p{0.68\linewidth}@{}}
    \toprule
    \textrm{\textbf{Parameter}} & \textbf{Try\ldots}\\\midrule
    sample, genes5, samples5 & other samples and genes in the $1\times1$ and $5\times5$ tables: which marker is high in which cell type?\\[2pt]
    mask & another cell type and cytokines: is the raw count higher with the cytokine?\\[2pt]
    MIN\_COUNT, MIN\_SAMPLES & a stricter or looser gene filter: how many genes survive?\\[2pt]
    CELL\_TYPE\_MARKERS, IMMUNE\_GENES & your own genes in the before/after heatmaps\\
    \bottomrule
  \end{tabular}
\end{frame}

\begin{frame}{Part 2 in one slide}
  \footnotesize
  \begin{tabular}{@{}c>{\raggedright\arraybackslash}p{0.14\linewidth}>{\raggedright\arraybackslash}p{0.33\linewidth}>{\raggedright\arraybackslash}p{0.22\linewidth}>{\raggedright\arraybackslash}p{0.17\linewidth}@{}}
    \toprule
    & \textbf{Step} & \textbf{Fixes} & \textbf{Formula} & \textbf{Stored in}\\\midrule
    1 & CPM & different library sizes (between samples) & $\text{count}_{ij}/L_i\times10^6$ & \code{layers["cpm"]}\\[2pt]
    2 & $\log_2(\cdot+1)$ & dynamic range, mean--variance link (across genes) & $\log_2(\text{CPM}+1)$ & \code{adata.X}\\[2pt]
    3 & gene filter & noise and too many tests & $\geq10$ counts in $\geq12$ samples & rows of \code{adata}\\[2pt]
    4 & z-score per gene & gene baseline hides the pattern (heatmaps) & $(x_{gi}-\bar x_g)/s_g$ & not stored\\
    \bottomrule
  \end{tabular}
  \vspace{0.4cm}
  \begin{block}{Remember}
    \code{layers["counts"]} = raw counts \quad $\to$ \quad \code{layers["cpm"]} = CPM \quad $\to$ \quad \code{adata.X} = log$_2$(CPM+1)
  \end{block}
\end{frame}

\begin{frame}[standout]
  Break\\[6pt]
  {\Large 10 minutes}
\end{frame}
